\documentclass[11pt,letterpaper]{article}

\usepackage[spanish,es-noshorthands]{babel}
\usepackage{fontspec}
\setmainfont{Source Serif Pro}
\usepackage{microtype}
\usepackage{geometry}
\usepackage{booktabs}
\usepackage{tabularx}
\usepackage{enumitem}
\usepackage{xcolor}
\usepackage{csquotes}
\usepackage{amsmath}
\usepackage{siunitx}
\usepackage{tikz}
\usetikzlibrary{arrows.meta,positioning}
\usepackage{xurl}
\usepackage[hidelinks]{hyperref}
\usepackage[backend=biber,style=apa,sorting=nyt]{biblatex}
\DeclareLanguageMapping{spanish}{spanish-apa}

\geometry{margin=2.3cm}
% Los rótulos de la figura usan nodos de ancho fijo con texto alineado a la
% derecha, y la bibliografía compone en modo sloppy: ambos generan avisos de
% línea holgada que son cosméticos y no indican un problema de composición.
% Se silencian para que `make check` siga siendo sensible a los desbordes reales.
\hbadness=10000
\setlength{\emergencystretch}{1em}
\setlength{\parindent}{0pt}
\setlength{\parskip}{0.55em}
\setlist{nosep}
\addbibresource{referencias.bib}
\AtBeginBibliography{\sloppy}

\sisetup{
  output-decimal-marker = {,},
  group-separator = {.},
  group-minimum-digits = 4,
  detect-all
}

% Colores de las dos generaciones químicas
\definecolor{cVieja}{RGB}{140,58,32}
\definecolor{cNueva}{RGB}{40,110,120}

% ---------------------------------------------------------------------------
% Figura de uso aparente: eje logarítmico de 10^2 a 10^7 kg de ingrediente
% activo. Posición horizontal: x = 1,9 * (log10(kg) - 2)
% ---------------------------------------------------------------------------
\newcommand{\molbar}[5]{%
  \node[left, text width=3.9cm, align=right, font=\scriptsize] at (-0.15,#1) {#4};%
  \fill[#3] (0,#1-0.15) rectangle (#2,#1+0.15);%
  \node[right, font=\scriptsize] at (#2+0.10,#1) {#5};%
}

\title{\vspace{-1.2cm}\textbf{Química de 1962 aplicada en 2022}\\
\large Por qué el registro de agroquímicos de Costa Rica mantiene viva
una tecnología obsoleta}
\author{Carlos Luis Rojas Aragonés\\
\small Gestión del Desarrollo Tecnológico}
\date{\small 15 de agosto de 2026}

\begin{document}

\maketitle
\vspace{-0.6em}

\section{Cómo la encontré}

Partí de un titular que se repite cada año en la prensa costarricense: según
FAOSTAT, Costa Rica encabeza el mundo en uso de plaguicidas por hectárea
\parencite{semanario2023}. El titular por sí solo no dice nada sobre tecnología,
así que fui a la fuente primaria, el informe de \emph{uso aparente} que publica
el Servicio Fitosanitario del Estado con las importaciones y exportaciones de
cada ingrediente activo \parencite{sfe2022}.

Al abrir el cuadro por molécula apareció el caso. No es que Costa Rica use
mucho plaguicida: es que \textbf{usa plaguicida antiguo}. De los 9,13 millones de
kilogramos de ingrediente activo de uso aparente en 2022, \textbf{el 70\,\%
corresponde a once moléculas comercializadas entre 1944 y 1974}. La química posterior
a 1995, que es la que la industria considera el estado del arte, aparece en el
mismo cuadro en cantidades de tres cifras. Es una tecnología teóricamente
jubilada que sigue siendo, en la práctica, la única que se aplica.

\section{La tecnología «obsoleta»}

Hablo de la primera generación de plaguicidas orgánicos de síntesis, la que se
desarrolló entre el final de la Segunda Guerra Mundial y los primeros años
setenta. Sus tres familias representativas en Costa Rica son:

\begin{itemize}
  \item \textbf{Ditiocarbamatos}, con el \textbf{mancozeb} (1962) como fungicida
    de contacto y acción multisitio.
  \item \textbf{Organofosforados y carbamatos}, con el \textbf{diazinón} (1952),
    el \textbf{clorpirifós} (1965), el \textbf{etoprofós} (1967) y el
    \textbf{oxamil} (1974), inhibidores de la acetilcolinesterasa.
  \item \textbf{Bipiridilos y fenilureas}, con el \textbf{paraquat} (1962) y el
    \textbf{diurón} (1954) como herbicidas.
\end{itemize}

Sus figuras de mérito \parencite{deweck2022} son las propias de un diseño
anterior a la biología molecular aplicada: actúan sobre procesos comunes a
muchos organismos, no sobre una diana concreta de la plaga; por eso hay que
aplicarlos en dosis de kilogramos por hectárea y por eso su toxicidad no
discrimina entre la plaga, el aplicador y el organismo acuático. El informe del
PNUD encuentra que 19 de los 21 plaguicidas más usados en el país, el 80\,\% del
volumen total, califican como altamente peligrosos según los criterios FAO/OMS,
y que de los 95 ingredientes activos asociados a efectos adversos para la salud
que se usan en Costa Rica, \textbf{56 no están autorizados en la Unión Europea}
\parencite{vargas2021}. Los dos fungicidas más representativos del bloque
perdieron allí su aprobación por vía reglamentaria: el clorotalonil en 2019 y el
mancozeb con vencimiento el 31 de enero de 2021
\parencite{ue2019677,ue20202087}.

\section{La tecnología llamada a reemplazarla}

La sustituta lleva treinta años disponible. Es la química de \textbf{sitio
único}, diseñada a partir del conocimiento del receptor concreto sobre el que
actúa: las diamidas antranílicas como el \textbf{clorantraniliprol}, que abre el
receptor de rianodina del insecto; los inhibidores de la biosíntesis de lípidos
como el \textbf{espirotetramato}; las espinosinas de origen fermentativo como el
\textbf{espinosad}; o las oxadiazinas como el \textbf{indoxacarb}. A su lado
corre el control biológico y la aplicación de precisión.

La mejora es de uno a dos órdenes de magnitud en la figura de mérito que importa,
la \textbf{dosis de ingrediente activo por hectárea necesaria para obtener
control}: entre \num{1500} y \num{2000} gramos por hectárea en el caso del
mancozeb frente a entre 10 y 50 gramos en el del clorantraniliprol, según las
dosis de etiqueta habituales. A eso se añade una toxicidad para mamíferos
sustancialmente menor y una persistencia ambiental más corta.

\section{El dato}

\begin{figure}[ht]
  \centering
  \begin{tikzpicture}
    % --- rejilla logarítmica ---
    \foreach \e in {2,3,4,5,6,7} {
      \pgfmathsetmacro{\xg}{1.9*(\e-2)}
      \draw[gray!22] (\xg,-0.30) -- (\xg,7.80);
      \node[below, font=\tiny, gray!55!black] at (\xg,-0.32)
        {$10^{\e}$};
    }
    \draw[gray!60!black] (0,-0.30) -- (9.5,-0.30);
    % --- generación anterior a 1975 ---
    \molbar{7.44}{8.590}{cVieja}{\textbf{Mancozeb} (1962)}{\num{3317217}}
    \molbar{7.00}{7.456}{cVieja}{2,4-D (1944)}{\num{840053}}
    \molbar{6.56}{7.194}{cVieja}{Diazinón (1952)}{\num{611566}}
    \molbar{6.12}{6.820}{cVieja}{Diurón (1954)}{\num{388649}}
    \molbar{5.68}{6.779}{cVieja}{Glifosato (1974)}{\num{369711}}
    \molbar{5.24}{6.334}{cVieja}{Oxamil (1974)}{\num{215538}}
    \molbar{4.80}{6.276}{cVieja}{Etoprofós (1967)}{\num{201096}}
    \molbar{4.36}{6.073}{cVieja}{Paraquat (1962)}{\num{157245}}
    \molbar{3.92}{6.019}{cVieja}{Clorpirifós (1965)}{\num{147258}}
    \molbar{3.48}{5.832}{cVieja}{Carbendazina (1972)}{\num{117367}}
    \molbar{3.04}{5.441}{cVieja}{Clorotalonil (1967)}{\num{73036}}
    % --- separador ---
    \draw[gray!45, densely dashed] (-4.15,2.71) -- (9.5,2.71);
    % --- generación posterior a 1995 ---
    \molbar{2.38}{3.655}{cNueva}{Espirotetramato (2007)}{\num{8385}}
    \molbar{1.94}{1.897}{cNueva}{Flupiradifurona (2015)}{997}
    \molbar{1.50}{1.367}{cNueva}{Espinosad (1997)}{524}
    \molbar{1.06}{1.365}{cNueva}{\textbf{Clorantraniliprol} (2008)}{523}
    \molbar{0.62}{1.037}{cNueva}{Acetamiprid (1995)}{351}
    \molbar{0.18}{0.874}{cNueva}{Indoxacarb (2000)}{288}
    % --- rótulos de bloque ---
    \node[right, font=\scriptsize\itshape, cVieja] at (6.7,3.04)
      {generación anterior a 1975};
    \node[right, font=\scriptsize\itshape, cNueva] at (5.9,1.94)
      {generación posterior a 1995};
    \node[below, font=\scriptsize, gray!45!black] at (4.75,-0.72)
      {Uso aparente en 2022, kilogramos de ingrediente activo (escala
       logarítmica)};
  \end{tikzpicture}
  \caption{Uso aparente de ingredientes activos en Costa Rica durante 2022,
  ordenado por volumen y separado por generación tecnológica. Elaboración propia
  con los datos del cuadro 2 del informe AE-REG-INF-001 del Servicio
  Fitosanitario del Estado \parencite{sfe2022}. Los años entre paréntesis
  corresponden a la primera comercialización o al primer registro de la molécula
  y tienen una incertidumbre de uno o dos años en varios casos; ninguna
  conclusión del artículo depende de esa precisión.}
  \label{fig:usoaparente}
\end{figure}

La Figura~\ref{fig:usoaparente} es el argumento entero en una imagen. El
fungicida más usado del país es una molécula de 1962 y absorbe por sí solo el
\textbf{36,3\,\%} de todo el ingrediente activo del año. La diamida moderna
equivalente, clorantraniliprol, registra 523 kilogramos. La razón entre ambas es
de \textbf{\num{6344} a 1} en masa.

Esa razón exagera la diferencia real de superficie tratada, porque la dosis por
hectárea de las dos moléculas no es comparable. Corregida por las dosis de
etiqueta señaladas antes, la relación cae a un orden de \textbf{50 a 100
hectáreas tratadas con química de 1962 por cada hectárea tratada con química de
2008}. Sigue siendo un país que, cuarenta y seis años después de que existiera
la alternativa, no la ha adoptado.

\section{Por qué no se ha abandonado}

\subsection*{01. La razón institucional: un registro asimétrico}

Esta es la causa dominante y la que hace el caso interesante para la gestión
tecnológica. En Costa Rica no existe una ley que prohíba importar moléculas
nuevas. Existe algo más eficaz: un régimen de registro que \textbf{no le pide
nada al producto viejo y se lo pide todo al producto nuevo}.

\begin{itemize}
  \item Los decretos ejecutivos 17557 (1986) y 24337 (1995) inscribieron
    plaguicidas \textbf{sin fecha de vencimiento} y con requisitos de información
    sobre peligrosidad que la propia Contraloría General de la República calificó
    de omisos. Sobre el segundo dictaminó que \enquote{se torna prácticamente
    inaplicable, por cuanto en muchas formas contradice, violenta, disminuye,
    sustituye o elimina los principios y normas establecidas en el ordenamiento
    jurídico} \parencite[citado en][]{vargas2021}.
  \item El decreto 33495 de 2006 introdujo por fin la exigencia del paquete
    completo de datos para el análisis de riesgo, pero \textbf{solo se aplica a
    los registros nuevos} \parencite{decreto33495}.
  \item La Ley 8702 de 2009 registró más de 400 plaguicidas con vigencia de diez
    años \textbf{sin evaluación ambiental ni sanitaria}, y otorgó al primer
    registrante diez años de protección de los datos de prueba, lo que cerró
    también la entrada de los genéricos \parencite{ley8702,vargas2021}.
\end{itemize}

Conviene precisar que el Estado sí actúa, pero solo en una dirección. En 2004 se
prohibieron por decreto 68 ingredientes activos y entre 2007 y 2008 se
publicaron diez decretos más de prohibición o restricción
\parencite{rojas2016}. El catálogo legal se encoge por un extremo sin que entre
nada por el otro, que es la forma más rápida de convertir un régimen de
protección en un régimen de congelación.

El resultado, según los oficios del propio Servicio Fitosanitario, es que
\textbf{\num{1884} plaguicidas formulados siguen comercializándose con el
registro vencido o sin fecha de vencimiento}, sin una segunda revisión ni de su
peligrosidad ni de su \emph{eficacia biológica}, pese a llevar quince, veinte,
treinta y más años en el mercado \parencite{vargas2021}. Del otro lado, el
Ministerio de Agricultura reconoció que el país pasó \textbf{veintiún años sin
poder registrar moléculas nuevas más eficientes} \parencite{delfino2024}.

\textcite{castrovargas2023} llaman a este equilibrio \emph{regulation by
impasse}: no es una decisión de política, es un atasco que se reproduce solo,
sostenido por la disputa entre ministerios, tribunales y autoridades
regulatorias, y que ninguna de las partes tiene incentivo suficiente para
resolver. Es exactamente el mecanismo de \emph{lock-in} descrito por
\textcite{arthur1989}, con la diferencia de que aquí el rendimiento creciente
que fija la trayectoria no es de mercado, sino administrativo.

La reforma llegó en diciembre de 2022 con el decreto 43838, que permite
homologar el registro ya otorgado en países de la OCDE \parencite{decreto43838}.
El tamaño de su efecto mide el tamaño del atasco previo: a octubre de 2024 se
habían registrado 39 productos, de los cuales 34 por homologación y
\textbf{solo 6 eran moléculas nuevas} \parencite{delfino2024}.

\subsection*{02. La razón económica: el genérico fuera de patente}

Las once moléculas del bloque superior de la figura están fuera de patente desde
hace décadas, se fabrican en varias plantas asiáticas y compiten solo por precio.
Las seis del bloque inferior son productos propietarios. El diferencial de coste
por hectárea tratada es de varias veces, y se compara contra un ingreso por
hectárea que no ha crecido en la misma proporción. Para el productor, el cálculo
es inmediato y correcto.

La oferta comercial refuerza el sesgo: el mancozeb acumula 839 registros de
producto formulado en el país y el paraquat 712 \parencite{vargas2021}. Cambiar
de molécula significa cambiar de proveedor, de distribuidor y de crédito de
insumos, no solo de envase.

\subsection*{03. La razón de conocimiento: 1 extensionista por cada 1.415
productores}

La química de sitio único exige saber qué plaga se tiene y en qué estadio, porque
fuera de su espectro simplemente no funciona. El producto de amplio espectro
perdona el error de diagnóstico. Y Costa Rica cuenta con unos 200 extensionistas
para unos \num{238000} productores, es decir un agente por cada \num{1415}
\parencite{delfino2023}. Cuando el servicio de extensión no existe, la tecnología
que sobrevive es la que tolera la ignorancia sobre la plaga.

\subsection*{04. La razón de riesgo percibido}

El coste de un fallo de control no es simétrico. Perder una cosecha de banano o
de piña por probar una molécula nueva es un daño concreto e inmediato; el daño
que causa la molécula vieja es difuso, diferido y lo paga en parte un tercero.
El país registró 58 muertes por intoxicación con plaguicidas entre 2010 y 2020 y
unos ₡\num{5000} millones anuales en costes sanitarios \parencite{delfino2023},
y el clorotalonil solo se prohibió en noviembre de 2023, después de que la
contaminación de nacientes obligara a abastecer a \num{10000} personas con
cisternas \parencite{delfino2023clorotalonil}.

\section{Por qué además ha renacido}

Hay una segunda respuesta, y es la que más me interesa como gestor de tecnología,
porque no tiene nada que ver con la inercia.

Las moléculas modernas deben su ventaja a que actúan sobre una sola diana. Esa
misma propiedad las hace frágiles: una única mutación en el patógeno anula el
producto. Las estrobilurinas, lanzadas a finales de los noventa como la
generación que jubilaba a los fungicidas de contacto, quedaron en buena medida
sin eficacia frente a \emph{Septoria} a comienzos de los años dos mil, y los
inhibidores SDHI tienen un riesgo de resistencia calificado de medio a alto
\parencite{ahdb2023}.

El mancozeb, en cambio, actúa sobre múltiples sitios a la vez. Está clasificado
como M03 en la lista del FRAC, el grupo de riesgo de resistencia bajo, y tras
seis décadas de uso intensivo los casos documentados de resistencia se limitan a
unas pocas poblaciones de \emph{Plasmopara viticola} en Europa
\parencite{gullino2010,frac2024}. Por eso la recomendación técnica actual no es
sustituirlo, sino \textbf{mezclarlo}: usar el multisitio de 1962 como socio de
mezcla que protege la vida útil de la molécula de sitio único
\parencite{ahdb2023}.

El atributo que hacía primitivo al mancozeb, su falta de selectividad, resultó
ser la ventaja decisiva en un eje que nadie estaba midiendo en 1990: la
\textbf{durabilidad de la eficacia}. El Cuadro~\ref{tab:fom} resume la
comparación completa.

\begin{table}[ht]
  \centering
  \small
  \caption{Comparación de las dos generaciones en las figuras de mérito
  relevantes.}
  \label{tab:fom}
  \begin{tabularx}{\textwidth}{@{}>{\raggedright\arraybackslash}p{0.24\textwidth}
    >{\raggedright\arraybackslash}X >{\raggedright\arraybackslash}X@{}}
    \toprule
    \textbf{Figura de mérito} & \textbf{Generación 1944--1974}
      & \textbf{Generación 1995--2015} \\
    \midrule
    Dosis para obtener control
      & \num{1500} a \num{2000} g ia/ha (mancozeb)
      & 10 a 50 g ia/ha (clorantraniliprol) \\
    Modo de acción
      & Multisitio, espectro amplio & Sitio único, selectivo \\
    Riesgo de resistencia
      & Bajo (FRAC M03) & Medio a alto \\
    Toxicidad y persistencia
      & 19 de los 21 más usados son altamente peligrosos (FAO/OMS)
      & Perfil toxicológico sustancialmente menor \\
    Estado en la Unión Europea
      & 56 de 95 ingredientes activos no autorizados
      & Autorizados \\
    Coste por hectárea
      & Genérico fuera de patente, varios oferentes
      & Producto propietario \\
    Exigencia de diagnóstico
      & Baja, tolera el error & Alta, exige identificar la plaga \\
    Estatus registral en Costa Rica
      & \num{1884} formulados sin vencimiento ni reválida
      & 21 años de bloqueo; 6 moléculas nuevas desde 2023 \\
    Uso aparente en 2022
      & 6,44 millones de kg ia (11 moléculas)
      & \num{11068} kg ia (6 moléculas) \\
    \bottomrule
  \end{tabularx}
\end{table}

\section{Precisiones y limitaciones}

\begin{enumerate}
  \item \textbf{Las cifras de kg/ha que circulan no son comparables entre sí.}
    El SFE calcula 8,84 kg ia/ha porque divide entre \num{1032320} hectáreas de
    uso agropecuario, pastos incluidos \parencite{sfe2022}; el informe del PNUD
    obtiene 34,45 kg ia/ha al dividir solo entre el área efectivamente cultivada
    \parencite{vargas2021}; FAOSTAT publica 23,4. Las tres son correctas con su
    denominador. Cito el dato oficial y señalo el desacuerdo en lugar de escoger
    el más llamativo.
  \item \textbf{El uso aparente no es uso.} Es importación menos exportación, sin
    considerar los saldos en bodega, y el informe de 2022 está marcado como
    preliminar \parencite{sfe2022}. Sirve para comparar órdenes de magnitud entre
    moléculas, no como medida de aplicación en campo.
  \item \textbf{La razón de \num{6344} a 1 es de masa, no de superficie.}
    Corregida por dosis de etiqueta cae a un orden de 50 a 100 a 1. Las dosis de
    etiqueta varían por cultivo y por plaga, de modo que esa corrección es una
    estimación de orden de magnitud, no un cálculo cerrado.
  \item \textbf{El glifosato y el 2,4-D siguen aprobados en la Unión Europea} y
    aparecen en el bloque antiguo solo por su año de introducción. La antigüedad
    de una molécula y su peligrosidad son dos ejes distintos, y el artículo
    sostiene el primero, no una equivalencia entre ambos.
  \item \textbf{El caso está en curso.} El decreto 43838 lleva tres años y medio
    vigente y su efecto todavía se mide en decenas de registros. Con los datos
    disponibles no puede afirmarse si la reforma desatasca el sistema o solo lo
    alivia.
\end{enumerate}

\section*{Lo que me llevo}

El caso obliga a corregir la intuición habitual sobre por qué sobrevive una
tecnología obsoleta. La respuesta cómoda es la inercia del usuario. Aquí la
respuesta es otra: durante veintiún años \textbf{el sistema regulatorio fue más
restrictivo con la tecnología mejor que con la peor}, porque la exigencia de
evidencia se aplicó a la entrada y no a la permanencia. Cualquier régimen que
evalúe lo nuevo con rigor y deje lo viejo sin caducidad produce el mismo
resultado, sea en agroquímicos, en dispositivos médicos o en sistemas
informáticos certificados.

Y una segunda lección, más incómoda: no todo lo que sobrevive lo hace por
inercia. El mancozeb también sigue ahí porque el eje en el que la generación
moderna es imbatible, la dosis por hectárea, deja fuera el eje en el que es
frágil, la durabilidad de la eficacia. Antes de declarar obsoleta una tecnología
conviene comprobar si el sucesor gana en todas las figuras de mérito o solo en
la que hoy resulta más visible.

\printbibliography[title={Referencias bibliográficas}]

\end{document}
